\documentclass[11pt,letterpaper]{article}

% Codificación e idioma
\usepackage[utf8]{inputenc}
\usepackage[T1]{fontenc}
\usepackage[spanish]{babel}

% Fuente Fourier (Utopia)
\usepackage{fourier}

% Fuente sans serif TeX Gyre Heros (solo para el encabezado)
\usepackage{tgheros}

% Márgenes de 2 cm
\usepackage[letterpaper,margin=2cm,headheight=14pt]{geometry}

% Encabezado y pie de página
\usepackage{fancyhdr}

% Textos del encabezado (se pueden cambiar en cualquier parte del documento)
\newcommand{\textoizquierdo}{Business Intelligence and Data Warehouse}
\newcommand{\textoderecho}{Module 1 - Introduction to Business Intelligence as an Analytical System}

% Uso en el documento:
%   \encabezadoizq{Nuevo texto}   -> cambia el texto de la izquierda
%   \encabezadoder{Module 2}     -> cambia el texto de la derecha
% El cambio aplica a la página donde se escribe; usa \clearpage antes
% si quieres que empiece en una página nueva.
\newcommand{\encabezadoizq}[1]{\renewcommand{\textoizquierdo}{#1}}
\newcommand{\encabezadoder}[1]{\renewcommand{\textoderecho}{#1}}

\pagestyle{fancy}
\fancyhf{}
\fancyhead[L]{\footnotesize\sffamily \textoizquierdo}
\fancyhead[R]{\footnotesize\sffamily\bfseries \textoderecho}
\fancyfoot[C]{\thepage}
\renewcommand{\headrulewidth}{0.4pt}
\renewcommand{\footrulewidth}{0pt}

% Plantilla para estilos de encabezado completos (descomentar y adaptar):
% \fancypagestyle{otroestilo}{%
%   \fancyhf{}
%   \fancyhead[L]{\footnotesize\sffamily Texto izquierdo}
%   \fancyhead[R]{\footnotesize\sffamily\bfseries Texto derecho}
%   \fancyfoot[C]{\thepage}}
% Activar con \pagestyle{otroestilo} o, para una sola página, \thispagestyle{otroestilo}.
\setlength{\parindent}{0pt}
 
\begin{document}

\section*{Introducción}
 
Este contenido introduce los conceptos fundamentales de Business Intelligence (BI) y su relación con los sistemas de data warehousing para apoyar la toma de decisiones empresariales.

\section{Introducción a Business Intelligence}
BI es una infraestructura tecnológica que mejora los procesos de negocio mediante 
la recopilación, integración, análisis y compartición de datos. Los sistemas comunes 
de BI incluyen sistemas de información ejecutiva, soporte a decisiones, procesamiento 
analítico en línea (OLAP) y gestión de relaciones con clientes. 
\section{Componentes y proceso de desarrollo de un sistema BI} 
El desarrollo comienza con la comprensión del negocio, 
seguido por la selección de fuentes de datos relevantes y el modelado multidimensional
(como esquemas estrella y copo de nieve). La integración de datos se realiza
mediante procesos ETL (extracción, transformación y carga), que abordan
problemas de calidad de datos. 
\section{Análisis y presentación de datos} 
El análisis básico se realiza con consultas SQL y OLAP para análisis descriptivo, mientras
que técnicas avanzadas incluyen minería de datos y análisis predictivo. La
información resultante debe presentarse de forma clara para facilitar la toma de
decisiones empresariales. 
\section{¿Cómo se analizan los datos?} 
Visualización de datos
\begin{itemize}
	\item Existen diversas formas de visualizar datos, como mapas digitales
(estáticos o animados), tablas, gráficos, realidad virtual y dashboards.
\end{itemize}


\begin{itemize}
	\item Estas visualizaciones permiten que los usuarios comprendan fácilmente
los resultados del análisis.
\end{itemize}
Componentes de Business Intelligence

\begin{itemize}
	\item Un sistema de soporte a decisiones involucra datos y modelos para ayudar
en la toma de decisiones.
\end{itemize}


\begin{itemize}
	\item Business Intelligence implica adquirir datos de múltiples fuentes,
consolidarlos en un data warehouse, analizarlos y presentar resultados visuales
para la toma de decisiones.
\end{itemize}
Técnicas de análisis de datos

\begin{itemize}
	\item El análisis descriptivo se realiza con consultas complejas y
multidimensionales, como consultas secuenciales y hojas de cálculo.
\end{itemize}


\begin{itemize}
	\item Técnicas más avanzadas como data mining permiten análisis predictivos y
prescriptivos, que serán detallados posteriormente.
\end{itemize}
Un data warehouse es un sistema de almacenamiento centralizado donde se
consolidan, integran y organizan datos provenientes de diversas fuentes. Su
función en Business Intelligence es: Reunir datos de diferentes sistemas para
tener una visión unificada. Facilitar el análisis mediante la estructuración y
limpieza de datos. Servir como base para realizar consultas, análisis
descriptivos, predictivos y prescriptivos. Apoyar la toma de decisiones al
proporcionar datos confiables y accesibles para visualizaciones y reportes.

El \textbf{modelado multidimensional} en Business Intelligence es una técnica
para organizar los datos en estructuras que facilitan el análisis desde
diferentes perspectivas o dimensiones.  \\ \\
Características clave:

\begin{itemize}
	\item Los datos se organizan en hechos (medidas cuantitativas, como ventas) y
dimensiones (categorías para analizar, como tiempo, ubicación, producto).
\end{itemize}


\begin{itemize}
	\item Permite realizar consultas complejas y análisis rápidos, como total de
ventas por ciudad o por período.
\end{itemize}


\begin{itemize}
	\item Facilita la creación de esquemas como Star Schema o Snowflake Schema,
que son modelos comunes en data warehouses.
\end{itemize}
Este modelado ayuda a que los usuarios puedan explorar y entender los datos de
forma intuitiva para tomar decisiones. 
\newpage
\encabezadoder{Module 2 - Design of a Data Warehouse}
\section{Diseñár e implementar una base de datos analítica}
Esta lección se centra en el diseño e implementación de una base de datos analítica 
para inteligencia de negocios, destacando el modelado multidimensional y la transición 
de sistemas transaccionales a sistemas OLAP. \\

\textbf{Estrategias de diseño de Data Warehouse}
	\begin{itemize}
		\item Existen dos enfoques principales: el de Inmon (top-down), que define un data warehouse corporativo primero, y el de Kimball (bottom-up), que construye data marts para necesidades específicas. 
	\end{itemize}
    
	\begin{itemize}
		\item El diseño se basa en un modelo multidimensional que incluye dimensiones, atributos, jerarquías y hechos o métricas.
	\end{itemize} 
    

\textbf{Componentes del modelo multidimensional}
	\begin{itemize}
		\item Las dimensiones representan áreas temáticas o líneas de negocio y contienen atributos organizados en jerarquías para facilitar el análisis. 
	\end{itemize}

	\begin{itemize}
		\item Los hechos son medidas o métricas que cuantifican el comportamiento del negocio y se almacenan en tablas centrales.
	\end{itemize}
    
    

\textbf{Proceso de modelado multidimensional}
	\begin{itemize}
		\item Primero se define un modelo de datos integrado para la empresa.
	\end{itemize}
    
	\begin{itemize}
		\item Luego se delimita el alcance del data warehouse, identificando qué partes del negocio y usuarios serán cubiertos.
    \end{itemize}

	\begin{itemize}
		\item Finalmente, se produce el modelo multidimensional que servirá de base para el esquema del data warehouse, permitiendo análisis detallados y agregados según diferentes dimensiones.
	\end{itemize}

\section{Pasos para el modelo multidimensional OLTP}
Hay 7 pasos:

\begin{enumerate}
	\item Definir dimensiones o las formas de analizar la información.
	\item Identificar los atributos de cada dimensión.
	\item Organizar jerárquicamente los atributos dentro de cada dimensión.
	\item Asignar identificadores únicos a los atributos.
	\item Asignar campos descriptivos en los atributos.
	\item Identificar los hechos y crear una tabla de hechos.
	\item Asignar identificadores únicos en la tabla de hechos.
\end{enumerate}

\end{document}
